\clearpage
\VerdierPhysicalPageAnchor{0220}%
\begingroup\small\setlength{\parskip}{0pt}\setlength{\emergencystretch}{1.5em}
\setlength{\abovedisplayskip}{4pt plus 2pt minus 2pt}%
\setlength{\belowdisplayskip}{4pt plus 2pt minus 2pt}%

\VerdierSemanticAnchor{ch3d-3-3-4-b}
\noindent\textit{b) Lorsque la catégorie $\mathcal C$ est additive, un
foncteur de $\mathcal A^{\operatorname{Ext}}$ dans $\mathcal C$ est
préadditif si et seulement si chacun des foncteurs de la suite de foncteurs
liés associée est additif.}

\VerdierSemanticAnchor{ch3d-3-3-4-proof}
Nous nous bornerons à montrer que $\Phi$ est essentiellement surjectif. Tout
d'abord, quitte à modifier les catégories en présence par une équivalence,
on peut supposer que la catégorie $\mathcal A$ possède un seul objet nul.
On vérifie aussitôt qu'alors la catégorie
$\mathcal A^{\operatorname{Ext}}$ ne possède qu'un seul objet nul. Soit
alors $(F^n,\delta_F^n)$ une suite de foncteurs liés. Remarquons que, pour
tout entier $n$, la propriété SFL II) définit un isomorphisme
$F^n(0)\longrightarrow F^{n+1}(0)$. On peut alors, en remplaçant au besoin
la suite $(F^n,\delta_F^n)$ par une suite isomorphe, supposer que, pour tout
$n$, $F^n(0)=F^{n+1}(0)$ et que
$\delta_F^n(0)=\operatorname{id}$.

\VerdierSemanticAnchor{ch3d-3-3-4-object-definition}
Tout objet non nul $W$ de la catégorie
$\mathcal A^{\operatorname{Ext}}$ s'écrit d'une manière unique sous la forme
$X[n]$, où $X$ est un objet de $\mathcal A$. Posons alors pour tout objet
non nul $W$ de $\mathcal A^{\operatorname{Ext}}$, $F(W)=F^n(X)$ et posons
de plus $F(0)=F^0(0)$. Nous avons ainsi défini une application de
$\operatorname{Ob}(\mathcal A^{\operatorname{Ext}})$ dans
$\operatorname{Ob}(\mathcal C)$ que nous allons maintenant prolonger en un
foncteur $F\colon\mathcal A^{\operatorname{Ext}}\longrightarrow\mathcal C$,
qui induise la suite de foncteurs liés $(F^n,\delta_F^n)$. Soit donc
$m\colon W\longrightarrow W'$ un morphisme de
$\mathcal A^{\operatorname{Ext}}$. Examinons d'abord différents cas :

\VerdierSemanticAnchor{ch3d-3-3-4-morphism-cases-zero}
\noindent 1) $W=W'=0$. On pose alors
$F(m)=\operatorname{id}_{F^0(0)}$.

\noindent 2) $W=0$ et $W'=X[p]$, $X\neq0$. On pose alors
$F(m)=F^p(0\longrightarrow X)$.

\noindent 3) $W=Y[p]$, $Y\neq0$, et $W'=0$. On pose alors
$F(m)=F^p(Y\longrightarrow0)$.

\VerdierSemanticAnchor{ch3d-3-3-4-degree}
Supposons maintenant que $W$ et $W'$ soient différents de l'objet nul. Ils
s'écrivent alors de manière unique sous la forme $W=X[p]$, $W'=Y[q]$, où
$X$ et $Y$ sont des objets de $\mathcal A$. L'entier $q-p$ est appelé alors
le \emph{degré} de $m$.

\VerdierSemanticAnchor{ch3d-3-3-4-morphism-cases-degree}
\noindent 4) Degré de $m<0$. On pose alors
$F(m)=F^q(0\longrightarrow Y)\,F^p(X\longrightarrow0)$.

\noindent 5) Degré de $m=0$. Le morphisme $m$ s'écrit alors de manière
unique sous la forme $f[p]$, où $f$ est un morphisme de $\mathcal A$. On
pose alors $F(m)=F^p(f)$.

\noindent 6) Degré de $m=1$. Il existe une suite exacte de $\mathcal A$ :
\[
S=0\longrightarrow Y\longrightarrow Z\longrightarrow X\longrightarrow0
\]
telle que $m=\delta^1(S)[p]$ (3.3.1). On pose alors
$F(m)=\delta_F^p(S)$. Le fait que le morphisme $F(m)$ ainsi défini ne
dépende pas de la suite exacte choisie résulte de (3.2.2) et de la propriété
SFL I).

\VerdierSemanticAnchor{ch3d-3-3-4-fl-le-n}
Désignons alors par
$\mathcal{F}\ell_{\leq n}(\mathcal A^{\operatorname{Ext}})$ l'ensemble des
morphismes $m\colon W\longrightarrow W'$ tels que $W$ ou $W'$ soit l'objet
nul ou tels que le degré de $m$ soit inférieur ou égal à $n$. L'existence du
foncteur $F\colon\mathcal A^{\operatorname{Ext}}\longrightarrow\mathcal C$
induisant la suite de foncteurs liés $(F^n,\delta_F^n)$ résulte alors du
lemme et de la proposition suivants :
\endgroup
